\section{Lý do chọn đề tài}
% TODO: Trình bày bối cảnh, vấn đề thực tiễn và lý do lựa chọn đề tài này.
% Gợi ý: Nêu xu hướng công nghệ, hạn chế của giải pháp hiện có, sự cần thiết của nghiên cứu.
\textit{(Nội dung sẽ được bổ sung)}

\section{Mục tiêu đề tài}
% TODO: Liệt kê các mục tiêu cụ thể, đo lường được mà đề tài hướng đến.
\begin{itemize}
    \item \textbf{Mục tiêu 1:} \textit{(Mô tả mục tiêu 1)}
    \item \textbf{Mục tiêu 2:} \textit{(Mô tả mục tiêu 2)}
    \item \textbf{Mục tiêu 3:} \textit{(Mô tả mục tiêu 3)}
\end{itemize}

\section{Phạm vi đề tài}
\subsection{Phạm vi về kỹ thuật}
% TODO: Nêu các công nghệ, nền tảng, công cụ được sử dụng trong đề tài.
\begin{itemize}
    \item \textit{(Công nghệ / Nền tảng 1)}
    \item \textit{(Công nghệ / Nền tảng 2)}
    \item \textit{(Công nghệ / Nền tảng 3)}
\end{itemize}

\subsection{Phạm vi về chức năng}
% TODO: Mô tả các chức năng / tính năng mà hệ thống sẽ triển khai.
\begin{itemize}
    \item \textit{(Chức năng 1)}
    \item \textit{(Chức năng 2)}
    \item \textit{(Chức năng 3)}
\end{itemize}

\section{Các đóng góp chính của đề tài}
% TODO: Tóm tắt những đóng góp mới, giá trị khoa học hoặc thực tiễn mà đề tài mang lại.
\begin{itemize}
    \item \textit{(Đóng góp 1)}
    \item \textit{(Đóng góp 2)}
    \item \textit{(Đóng góp 3)}
\end{itemize}

\section{Cấu trúc Đồ án cuối kì}
% TODO: Mô tả ngắn gọn nội dung từng chương trong đồ án.
\begin{itemize}
    \item \textbf{Chương 1. Giới thiệu đề tài:} \textit{(Tóm tắt nội dung chương 1)}
    \item \textbf{Chương 2. Cơ sở lý thuyết:} \textit{(Tóm tắt nội dung chương 2)}
    \item \textbf{Chương 3. Phương pháp luận:} \textit{(Tóm tắt nội dung chương 3)}
    \item \textbf{Chương 4. Triển khai và kiểm thử:} \textit{(Tóm tắt nội dung chương 4)}
    \item \textbf{Chương 5. Đánh giá và hướng phát triển:} \textit{(Tóm tắt nội dung chương 5)}
\end{itemize}
